% !TEX program = xelatex
% Bien dich: xelatex finsight_spec_chitiet.tex (chay 2 lan de cap nhat muc luc)
\documentclass[11pt,a4paper]{article}

\usepackage{fontspec}
\usepackage{amsmath}
\usepackage[a4paper,margin=2.2cm]{geometry}
\usepackage{enumitem}
\usepackage{booktabs}
\usepackage{longtable}
\usepackage{array}
\usepackage{xcolor}
\usepackage{tcolorbox}
\usepackage{fancyhdr}
\usepackage{titlesec}
\usepackage[hidelinks]{hyperref}

\definecolor{fsblue}{RGB}{20,80,140}
\definecolor{fsgreen}{RGB}{20,120,70}
\definecolor{fsgray}{RGB}{90,90,90}
\definecolor{fslight}{RGB}{238,243,250}
\definecolor{codebg}{RGB}{246,246,246}

\hypersetup{colorlinks=true, linkcolor=fsblue, urlcolor=fsblue}

\titleformat{\section}{\Large\bfseries\color{fsblue}}{\thesection}{0.6em}{}
\titleformat{\subsection}{\large\bfseries\color{fsblue!85!black}}{\thesubsection}{0.6em}{}
\titleformat{\subsubsection}{\normalsize\bfseries\color{fsgreen!75!black}}{\thesubsubsection}{0.5em}{}

\setlist[itemize]{leftmargin=1.35em, itemsep=2pt, topsep=2pt}
\setlist[enumerate]{leftmargin=1.55em, itemsep=2pt, topsep=2pt}

% Khoi code: ASCII-only ben trong de an toan voi xelatex
\usepackage{fvextra}
\fvset{frame=single, framesep=4pt, fontsize=\small, breaklines=true,
  rulecolor=\color{fsblue!40}}

\newtcolorbox{ghichu}{colback=fslight, colframe=fsblue, boxrule=0.6pt, arc=2pt,
  left=6pt, right=6pt, top=4pt, bottom=4pt}
\newtcolorbox{quyetdinh}{colback=fsgreen!7, colframe=fsgreen!70!black, boxrule=0.6pt,
  arc=2pt, left=6pt, right=6pt, top=4pt, bottom=4pt}

\newcommand{\dec}[1]{\textbf{\color{fsgreen!60!black}[QUYẾT ĐỊNH]} #1}

\pagestyle{fancy}
\fancyhf{}
\lhead{\small\color{fsgray} FinSight --- Spec chi tiết}
\rhead{\small\color{fsgray} v1}
\cfoot{\small\thepage}
\renewcommand{\headrulewidth}{0.3pt}

\title{\vspace{-1.2cm}\textbf{\color{fsblue}FinSight --- Bản đặc tả kỹ thuật chi tiết}\\[4pt]
\large Implementation Spec theo từng pha\\[5pt]
\normalsize\color{fsgray} Mỗi quyết định: thư viện, tham số, schema, thuật toán, fallback}
\author{}
\date{}

\begin{document}
\maketitle
\thispagestyle{fancy}
\vspace{-0.9cm}

\begin{ghichu}
\textbf{Mục đích tài liệu.} Đây là spec triển khai --- trả lời ``làm \emph{như thế nào} cụ thể'', không phải ``làm \emph{cái gì}''. Mọi lựa chọn đều kèm tham số mặc định, ngưỡng, và đường fallback. Tài liệu này đi cùng \texttt{finsight\_ke\_hoach.tex} (bản tổng quan kiến trúc).
\end{ghichu}

\tableofcontents
\vspace{0.3cm}

% =====================================================================
\section{Quy ước chung}

\begin{itemize}
  \item \textbf{Ngôn ngữ:} Python 3.11; quản lý phụ thuộc bằng \texttt{uv} (hoặc Poetry).
  \item \textbf{Đơn vị tiền chuẩn hóa nội bộ:} VND (đồng), kiểu \texttt{int64} hoặc \texttt{Decimal} --- KHÔNG dùng float cho tiền.
  \item \textbf{Định danh:} \texttt{doc\_id} (1 file BCTC), \texttt{report\_id} (1 báo cáo trong file: BS/IS/CF/Notes), \texttt{nam} (năm tài chính), \texttt{ma\_so} (mã số chỉ tiêu TT200).
  \item \textbf{Provenance bắt buộc:} mọi số liệu mang theo \texttt{(doc\_id, page, bbox)} để truy vết \& trích dẫn.
\end{itemize}

% =====================================================================
\section{Pha 1 --- OCR giữ cấu trúc (chi tiết)}

\subsection{Phân loại trang: digital vs scan}
\begin{quyetdinh}
\dec{} PyMuPDF (\texttt{fitz}). Với mỗi trang, đếm ký tự text trích được. Ngưỡng: \texttt{n\_chars >= 50} \emph{và} diện tích text-span phủ $> 5\%$ trang $\Rightarrow$ trang \emph{digital}; ngược lại $\Rightarrow$ \emph{scan}.
\end{quyetdinh}
\begin{itemize}
  \item Digital: lấy text + bbox bằng \texttt{page.get\_text("dict")} (blocks $\rightarrow$ lines $\rightarrow$ spans, có bbox). Bảng thử trước bằng \texttt{page.find\_tables()}.
  \item Scan: render ảnh \texttt{page.get\_pixmap(dpi=300)} $\rightarrow$ đưa sang pipeline OCR.
  \item Trang lai (vừa text vừa ảnh chèn): xử lý theo vùng --- vùng ảnh OCR, vùng text đọc thẳng.
\end{itemize}

\subsection{Pipeline OCR + layout cho trang scan}
\begin{quyetdinh}
\dec{} Bố cục \& cấu trúc bảng: \textbf{PaddleOCR PP-Structure} (model layout \texttt{PP-DocLayout}; table structure \texttt{SLANet}). Nhận dạng \emph{chữ tiếng Việt}: \textbf{VietOCR} (\texttt{vgg\_transformer}) chạy trên từng dòng cắt ra --- vì recognition đa ngữ của Paddle yếu với dấu tiếng Việt.
\end{quyetdinh}

Pipeline từng bước:
\begin{enumerate}
  \item PP-DocLayout $\rightarrow$ phân vùng: \{title, text, table, figure, header, footer, page-number\}.
  \item Với vùng \texttt{text}: PP-Structure text-detection $\rightarrow$ box từng dòng $\rightarrow$ crop $\rightarrow$ \textbf{VietOCR} nhận dạng.
  \item Với vùng \texttt{table}: SLANet $\rightarrow$ cấu trúc bảng (HTML + bbox từng ô + row/col span) $\rightarrow$ crop từng ô $\rightarrow$ VietOCR nhận dạng nội dung ô.
  \item Thứ tự đọc (reading order): thuật toán \textbf{XY-cut} đệ quy; nếu 1 cột thì sort theo \texttt{top} rồi \texttt{left}.
\end{enumerate}

\subsection{Thuật toán hậu xử lý ô gộp \& phân cấp hàng (BCTC)}
Bảng cân đối có \emph{thụt lề mang ngữ nghĩa} (Tài sản ngắn hạn $>$ Tiền $>$ Tiền mặt). Thuật toán gán cấp:
\begin{enumerate}
  \item Lấy tọa độ \texttt{x\_left} của text trong ô nhãn (cột đầu) của mọi hàng.
  \item Gom cụm các giá trị \texttt{x\_left} thành các \emph{mức thụt lề} (làm tròn về bội số $\approx$12--16px, hoặc 1D k-means k=3..5).
  \item Hàng có mức thụt sâu hơn = con của hàng gần nhất phía trên có mức thụt nông hơn $\Rightarrow$ dựng cây phân cấp.
  \item Ô gộp (row/col span từ SLANet): nhân bản giá trị cho các ô con khi cần, đánh dấu \texttt{is\_merged}.
  \item Gán cột số $\rightarrow$ kỳ: đọc hàng header (``Số cuối năm'' / ``Số đầu năm'') để map cột $\rightarrow$ \texttt{nam}.
\end{enumerate}

\subsection{Schema đầu ra (document model)}
\begin{Verbatim}
{
  "doc_id": "vnm_2023",
  "pages": [
    {
      "page": 4, "is_scanned": true,
      "blocks": [
        {"type":"title","text":"BANG CAN DOI KE TOAN","bbox":[..]},
        {"type":"table","table_id":"t1",
         "cells":[
           {"row":3,"col":0,"text":"Tien va tuong duong tien",
            "indent_level":1,"row_span":1,"col_span":1,"bbox":[..]},
           {"row":3,"col":1,"text":"110","bbox":[..]},
           {"row":3,"col":2,"text":"12.345.678.901","bbox":[..]}
         ],
         "col_header_map":{"2":"so_cuoi_nam","3":"so_dau_nam"}}
      ]
    }
  ]
}
\end{Verbatim}

\subsection{Đánh giá Pha 1}
\begin{itemize}
  \item Text: \textbf{CER}/\textbf{WER} bằng \texttt{jiwer}. Mục tiêu CER $<$ 3\% trên trang digital, $<$ 8\% trên scan.
  \item Cấu trúc bảng: \textbf{TEDS} (Tree-Edit-Distance Similarity, script gốc PubTabNet). Mục tiêu TEDS $>$ 0.90.
\end{itemize}

% =====================================================================
\section{Pha 2 --- Trích xuất \& Chuẩn hóa (chi tiết)}

\subsection{Mô hình trích xuất bảng (phần DL lõi)}
\begin{quyetdinh}
\dec{} Hai phương án, chọn theo dữ liệu sẵn có:
\begin{itemize}
  \item \textbf{A --- Decoupled (khuyên dùng cho MVP):} Table Transformer (TATR, DETR-based) cho \emph{structure recognition} + VietOCR cho nội dung ô. Dễ debug, ít dữ liệu, lỗi tách bạch.
  \item \textbf{B --- End-to-end:} \texttt{Qwen2-VL-2B-Instruct} (hoặc Donut) finetune ảnh-bảng $\rightarrow$ JSON. Mạnh hơn khi bảng phức tạp, nhưng cần nhiều dữ liệu \& khó debug.
\end{itemize}
MVP: đi phương án A; chuẩn bị dữ liệu để nâng cấp B sau.
\end{quyetdinh}

\subsubsection{Cấu hình LoRA (khi finetune phương án B --- Qwen2-VL-2B)}
\begin{Verbatim}
LoraConfig(
  r=16, lora_alpha=32, lora_dropout=0.05, bias="none",
  task_type="CAUSAL_LM",
  target_modules=["q_proj","k_proj","v_proj","o_proj",
                  "gate_proj","up_proj","down_proj"],
  # tuy chon: them lop merger cua vision encoder
)
# Train: lr=1e-4, cosine schedule, warmup_ratio=0.03,
#        bf16, grad_accum -> effective batch=16, epochs=3..5,
#        max_grad_norm=1.0. Framework: PEFT + transformers/TRL.
\end{Verbatim}
Donut (encoder-decoder MBART): LoRA \texttt{r=8..16}, \texttt{target\_modules=["q\_proj","k\_proj","v\_proj","out\_proj"]} ở decoder.

\subsubsection{Xây dữ liệu huấn luyện}
\begin{itemize}
  \item \textbf{Annotation:} \texttt{PPOCRLabel} (xuất cấu trúc bảng HTML, hợp PaddleOCR) hoặc \texttt{Label Studio} (template OCR/table). Nhãn: bbox ô + text + row/col span + mã số.
  \item \textbf{Tổng hợp (synthetic) để pretrain:} render hàng nghìn ảnh bảng BCTC từ \emph{chart of accounts TT200} + số ngẫu nhiên hợp lệ, đa font/độ nhiễu/đường kẻ. Ground-truth sinh kèm $\Rightarrow$ miễn phí, vô hạn.
  \item \textbf{Thật:} 15--20 bảng làm golden eval; $\sim$50--200 bảng người sửa (từ HITL) để finetune.
  \item \textbf{Chiến lược:} pretrain trên synthetic $\rightarrow$ finetune trên dữ liệu thật đã sửa.
  \item \textbf{Target format (phương án B):} chuỗi JSON các hàng/ô (xem schema dưới).
\end{itemize}

\subsubsection{Pipeline inference từng bước}
\begin{enumerate}
  \item Ảnh trang $\rightarrow$ phát hiện vùng bảng (PP-DocLayout).
  \item Crop bảng $\rightarrow$ TATR (structure) \emph{hoặc} Qwen2-VL (JSON).
  \item (Phương án A) Crop từng ô $\rightarrow$ VietOCR $\rightarrow$ text.
  \item Lắp lưới ô + áp thuật toán phân cấp/ô gộp (Pha 1).
  \item \textbf{Parse số} (xem dưới) + phát hiện đơn vị tiền.
  \item \textbf{Chuẩn hóa nhãn $\rightarrow$ mã số TT200} (mục \ref{sec:tt200}).
  \item \textbf{Validation số học} (mục \ref{sec:valid}).
\end{enumerate}

\subsection{Parse số tiếng Việt}
\begin{quyetdinh}
\dec{} Regex + luật: dấu chấm là phân tách nghìn (\texttt{12.345.678}); số âm trong ngoặc \texttt{(1.234)} $\rightarrow$ \texttt{-1234}; ô \{``-'', ``'', ``N/A''\} $\rightarrow$ \texttt{null}; nhân hệ số đơn vị (``triệu đồng'' $\rightarrow \times 10^6$). Lưu dạng \texttt{int} VND.
\end{quyetdinh}

\subsection{Chuẩn hóa $\rightarrow$ mã số TT200 (cascade 4 tầng)}
\label{sec:tt200}
\begin{ghichu}
\textbf{Insight cốt lõi:} BCTC theo TT200 \emph{đã in sẵn cột ``Mã số''} (110, 200, 311\dots). Vì vậy đường chính KHÔNG phải fuzzy match nhãn, mà là \textbf{đọc thẳng cột mã số} rồi tra bảng tĩnh. Fuzzy/embedding/LLM chỉ là fallback khi mã số mất/mờ.
\end{ghichu}

\begin{enumerate}
  \item \textbf{Tầng 1 --- Đọc mã số trực tiếp (chính, độ chính xác cao nhất):} lấy giá trị cột ``Mã số'' $\rightarrow$ tra \texttt{TT200\_CHART} (dict tĩnh $\sim$ vài trăm mã, seed từ phụ lục Thông tư 200). Hợp lệ $\Rightarrow$ chấp nhận.
  \item \textbf{Tầng 2 --- Fuzzy match nhãn} (khi mã số thiếu/garbled): thư viện \textbf{RapidFuzz}, hàm \texttt{token\_sort\_ratio}, \textbf{ngưỡng $\geq 88$}. Tiền xử lý: lowercase, bỏ dấu câu, chuẩn hóa khoảng trắng. Khớp $\Rightarrow$ chấp nhận.
  \item \textbf{Tầng 3 --- Embedding similarity} (khi fuzzy $<88$): model \texttt{bkai-foundation-models/vietnamese-bi-encoder} (hoặc \texttt{bge-m3}); \textbf{cosine}; encode nhãn vs embedding của toàn bộ nhãn chuẩn TT200; lấy top-1, \textbf{chấp nhận nếu cos $> 0.75$}.
  \item \textbf{Tầng 4 --- LLM classify} (khi cos $\leq 0.75$ hoặc top-2 sát nhau): đưa nhãn + top-5 ứng viên (từ tầng 3) cho LLM chọn, few-shot 3--5 ví dụ, \emph{output ràng buộc} (chỉ 1 mã hoặc \texttt{UNKNOWN}).
  \item \textbf{Fallback cuối:} không tầng nào đạt $\Rightarrow$ gán \texttt{unmapped}, lưu nhãn thô, \textbf{đẩy HITL}. Tuyệt đối không bỏ âm thầm.
\end{enumerate}

Prompt LLM (tầng 4):
\begin{Verbatim}
He thong: Ban la chuyen gia ke toan VN (TT200). Cho mot ten chi tieu
va danh sach ma so ung vien, chon DUNG MOT ma phu hop nhat, hoac tra
"UNKNOWN" neu khong chac. Chi tra ma so.

Vi du:
- "Tien va cac khoan tuong duong tien" -> 110
- "Nguoi mua tra tien truoc ngan han" -> 312
- "Chi phi tra truoc dai han" -> 261

Ten chi tieu: "{label}"
Ung vien: {top5_codes_with_names}
Tra loi:
\end{Verbatim}

\textbf{Từ điển đồng nghĩa:} seed thủ công từ phụ lục TT200; \emph{tự lớn lên} --- mỗi correction của người ở tầng fallback ghi thêm 1 cặp (nhãn thô $\rightarrow$ mã chuẩn) vào dictionary, lần sau khớp ngay ở tầng 2.

\subsection{Validation số học (bắt lỗi extraction)}
\label{sec:valid}
\begin{itemize}
  \item \texttt{Tong tai san (270) == Tai san NH (100) + Tai san DH (200)}.
  \item \texttt{Tong cong nguon von (440) == Tong tai san (270)}.
  \item Cha = tổng các con theo cây phân cấp (sai số cho phép do làm tròn $\leq$ 1 đơn vị).
  \item Vi phạm $\Rightarrow$ cờ \texttt{extraction\_warning} + đẩy HITL ô liên quan.
\end{itemize}

\subsection{Schema chuẩn hóa (1 dòng, lưu Postgres)}
\begin{Verbatim}
{ "doc_id":"vnm_2023","report":"balance_sheet","nam":2023,
  "ma_so":"110","chi_tieu":"Tien va tuong duong tien",
  "parent_ma_so":"100","indent_level":1,"thuyet_minh":"V.01",
  "gia_tri": 12345678901, "don_vi":"VND",
  "map_method":"direct_code","map_confidence":1.0,
  "provenance": {"page":4,"bbox":[..]} }
\end{Verbatim}

% =====================================================================
\section{Pha 3 --- GraphRAG đa năm (chi tiết)}

\begin{quyetdinh}
\dec{} \textbf{Neo4j} cho đồ thị; \textbf{Microsoft GraphRAG} (hoặc pipeline tự xây) cho indexing + community summary; embedding \texttt{bge-m3} cho truy hồi ngữ nghĩa.
\end{quyetdinh}

Mô hình đồ thị (Cypher rút gọn):
\begin{Verbatim}
(:Company {mst, ten})
(:LineItem {ma_so, chi_tieu})
(:Value {gia_tri, nam, report})
(:Note {ma, noi_dung})

(LineItem)-[:HAS_VALUE]->(Value)
(Value)-[:OF_YEAR]->(:Year {nam})
(LineItem)-[:CHILD_OF]->(LineItem)        // phan cap
(Value)-[:YOY {delta_pct}]->(Value)        // nam-qua-nam
(LineItem)-[:REFERENCES]->(Note)           // tham chieu thuyet minh
(LineItem)-[:COMPUTED_FROM]->(LineItem)    // vi du LN gop = DT - GV
\end{Verbatim}
\begin{itemize}
  \item Cạnh \texttt{YOY} tính sẵn lúc ingest $\Rightarrow$ truy vấn xu hướng nhanh.
  \item Community summary: gom cụm theo nhóm chỉ tiêu (thanh khoản, đòn bẩy\dots) để trả lời câu hỏi ``toàn cục''.
  \item Chống bùng nổ node: chỉ mô hình hóa chỉ tiêu cốt lõi + thuyết minh quan trọng (whitelist mã số).
\end{itemize}

% =====================================================================
\section{Pha 4 --- Conformal Prediction (chi tiết)}

Áp dụng cho \emph{bài toán phân loại ánh xạ nhãn $\rightarrow$ mã số TT200} (và tương tự cho cờ rủi ro).

\begin{quyetdinh}
\dec{} \textbf{Split Conformal} với \textbf{APS} (Adaptive Prediction Sets). Coverage target $1-\alpha = 0.90$.
\end{quyetdinh}

\begin{itemize}
  \item \textbf{Điểm số:} bộ phân loại (embedding/LLM) cho phân phối $\hat p(y\mid x)$ trên các mã ứng viên.
  \item \textbf{Nonconformity score (APS):} với mẫu calibration $(x,y)$, sắp xếp $\hat p$ giảm dần, cộng dồn đến khi gặp lớp đúng $y$ $\Rightarrow$ $s = $ tổng tích lũy đó.
  \item \textbf{Calibration set:} 200--500 mẫu \emph{đã có nhãn người} (golden + correction), KHÔNG dùng để train.
  \item \textbf{Ngưỡng:} $\hat q = $ phân vị $\lceil (n{+}1)(1-\alpha)\rceil / n$ của tập điểm calibration.
  \item \textbf{Prediction set:} $C(x) = \{y : \text{APS score}(x,y) \le \hat q\}$.
  \item \textbf{Luật abstain $\rightarrow$ HITL:}
  \begin{itemize}
    \item $|C(x)| = 1 \Rightarrow$ tự động chấp nhận.
    \item $|C(x)| > 1$ (mơ hồ) hoặc $|C(x)| = 0$ (lạ) $\Rightarrow$ \textbf{abstain}, đẩy người duyệt.
  \end{itemize}
  \item \textbf{Khi có drift:} theo dõi \emph{độ phủ thực nghiệm} trên cửa sổ trượt; nếu tụt dưới $1-\alpha \Rightarrow$ recalibrate, hoặc dùng \textbf{Adaptive Conformal Inference (ACI)} cập nhật $\alpha_t$ online.
  \item \textbf{Thư viện:} \texttt{MAPIE} (hoặc tự cài $\sim$30 dòng).
\end{itemize}

% =====================================================================
\section{Pha 5 --- LangGraph Agent (chi tiết)}

\subsection{Graph state schema}
\begin{Verbatim}
class AgentState(TypedDict):
    session_id: str
    task: str                 # nhiem vu nguoi dung
    plan: list[str]           # buoc da phan ra
    cursor: int               # buoc hien tai
    messages: list            # hoi thoai (LangChain msgs)
    retrieved: list[dict]     # so lieu lay tu tool (co provenance)
    citations: list[dict]     # nguon trich dan
    confidence: float
    needs_verification: bool
    answer: str | None
    abstained: bool
\end{Verbatim}

\subsection{Các node \& trách nhiệm}
\begin{longtable}{>{\raggedright\arraybackslash}p{3.0cm} >{\raggedright\arraybackslash}p{12cm}}
\toprule
\textbf{Node} & \textbf{Làm gì} \\
\midrule
\endhead
\texttt{planner} & Phân rã \texttt{task} thành \texttt{plan} (chọn loại nhiệm vụ + tool cần dùng). \\
\texttt{router} & Điều phối: extract / compare / reconcile / qa; chọn tool kế tiếp theo \texttt{cursor}. \\
\texttt{tool\_executor} & Gọi tool, ghi kết quả vào \texttt{retrieved} kèm provenance. \\
\texttt{verifier} & Nếu \texttt{needs\_verification}: RAG quy định + web; tính \texttt{confidence}. \\
\texttt{synthesizer} & Viết \texttt{answer} + gắn \texttt{citations}. \\
\texttt{abstain} & Trả ``không đủ cơ sở'' khi thiếu bằng chứng / confidence thấp. \\
\bottomrule
\end{longtable}

\subsection{Cạnh \& điều kiện (conditional edges)}
\begin{Verbatim}
planner -> router
router  -> tool_executor
tool_executor -> router        # neu con buoc: cursor < len(plan)
tool_executor -> verifier      # neu het buoc
verifier -> synthesizer        # if confidence >= 0.7 and citations != []
verifier -> tool_executor      # if thieu du lieu (can lay them)
verifier -> abstain            # if confidence < 0.5 or no evidence
synthesizer -> END
abstain -> END
\end{Verbatim}
Hàm điều kiện ví dụ: \texttt{should\_verify(state)} trả \texttt{"verify"} nếu task chạm quy định/cờ đỏ hoặc \texttt{confidence < 0.7}.

\subsection{Tools (chữ ký + luật)}
\begin{Verbatim}
query_graph(intent|cypher) -> rows[+provenance]
get_lineitem(ma_so, nam)   -> {gia_tri, provenance}
compute_ratio(name, nam)   -> {value, formula, inputs[+provenance]}
reconcile(scope)           -> [discrepancy{a,b,delta,severity}]
search_regulation(query)   -> [chunk{text, cite}]
web_search(query)          -> [result{title,url,snippet}]
\end{Verbatim}
\begin{ghichu}
\textbf{Luật chống ảo giác (bắt buộc):} agent CHỈ được lấy con số qua tool; cấm tự ``nhớ''/sinh số. Mọi số trong \texttt{answer} phải truy được về \texttt{provenance} của một tool-call. Synthesizer từ chối câu trả lời chứa số không có nguồn.
\end{ghichu}

\subsection{Memory store}
\begin{itemize}
  \item \textbf{Ngắn hạn:} \texttt{messages} trong state; khi $>$ N token $\Rightarrow$ tóm tắt (summary node) giữ ngữ cảnh.
  \item \textbf{Dài hạn (theo session = 1 hồ sơ DN):}
  \begin{itemize}
    \item \textbf{Qdrant} 1 collection, record \texttt{\{session\_id, type:finding|decision|fact, content, embedding, citation, approved\_by, ts\}}; truy hồi top-k lọc theo \texttt{session\_id}.
    \item \textbf{Postgres} bảng \texttt{approved\_findings} cho recall chính xác + audit.
  \end{itemize}
  \item \textbf{Checkpoint graph:} \texttt{PostgresSaver} (LangGraph) lưu state $\Rightarrow$ mở lại session, resume đúng chỗ.
\end{itemize}

% =====================================================================
\section{Pha 6 --- Vòng HITL \& Finetune (chi tiết)}

\subsection{Định dạng dữ liệu correction (JSONL)}
\begin{Verbatim}
{ "doc_id":"vnm_2023","report":"balance_sheet","page":4,
  "table_id":"t1","cell":{"row":3,"col":2},
  "field":"gia_tri",            // gia_tri | ma_so | chi_tieu
  "predicted":"12.345.678.001","corrected":"12.345.678.901",
  "image_crop_path":"crops/vnm_2023_p4_t1_r3c2.png",
  "map_method_at_pred":"direct_code","confidence_at_pred":0.82,
  "corrected_by":"analyst_03","ts":"2026-06-19T10:00:00Z" }
\end{Verbatim}

\subsection{Trigger huấn luyện lại (drift-based, không theo lịch)}
Kích hoạt khi MỘT trong các điều kiện:
\begin{itemize}
  \item Độ phủ conformal thực nghiệm (cửa sổ trượt) $< 1-\alpha$ (vd $<0.90$).
  \item PSI của phân phối confidence $> 0.2$ so với baseline.
  \item Số correction tích lũy $\geq 50$.
\end{itemize}

\subsection{Quy trình finetune}
\begin{enumerate}
  \item Dựng tập train mới = synthetic gốc + toàn bộ correction (\textbf{oversample correction $\times 3$}).
  \item Train \textbf{adapter LoRA MỚI} (\texttt{adapter\_vN+1}), KHÔNG ghi đè adapter đang chạy. Cấu hình như Pha 2.
  \item Đánh giá trên \textbf{golden eval set ĐÓNG BĂNG} (không đổi giữa các vòng).
  \item Đăng ký adapter + metrics vào \textbf{MLflow}.
\end{enumerate}

\subsection{Chiến lược merge \& phục vụ}
\begin{quyetdinh}
\dec{} Không merge vào base mặc định. Phục vụ \textbf{base + active adapter} qua \textbf{vLLM} (hot-swap LoRA). Chỉ \texttt{merge\_and\_unload} khi ra bản ổn định để giảm overhead inference. Mỗi adapter versioned, ``active'' trỏ bằng con trỏ \texttt{active\_adapter\_id} trong DB/config.
\end{quyetdinh}

\subsection{Cổng promote \& rollback}
\begin{itemize}
  \item \textbf{Promote} adapter mới $\rightarrow$ active CHỈ khi: accuracy golden $\geq$ active hiện tại $- \varepsilon$ (\textbf{$\varepsilon=0.5\%$}) \emph{và} không regression ở bất kỳ nhóm chỉ tiêu nào (per-category breakdown).
  \item \textbf{Shadow trước promote:} chạy adapter mới ở chế độ shadow trên mẫu, so sánh, rồi mới đổi con trỏ.
  \item \textbf{Rollback} = lật \texttt{active\_adapter\_id} về bản trước (O(1), không tải lại base). Giữ tối thiểu 3 bản gần nhất.
  \item Ghi log + alert khi regression.
\end{itemize}

% =====================================================================
\section{Pha 0/7/8/9 --- Hạ tầng, hiệu năng, đánh giá, vận hành (chi tiết ngắn)}

\subsection{Pha 0 --- Hạ tầng}
\texttt{docker-compose}: \texttt{postgres(+pgvector)}, \texttt{redis}, \texttt{minio}, \texttt{qdrant}, \texttt{neo4j}. Quản version dữ liệu/model bằng \textbf{DVC} + \textbf{MLflow}.

\subsection{Pha 7 --- Hiệu năng \& song song}
\begin{itemize}
  \item \textbf{Celery + Redis}: 1 file = 1 task; concurrency = (số core $-2$). OCR trang song song trong file.
  \item \textbf{Idempotency:} key = SHA-256 file; đã xử lý $\Rightarrow$ bỏ qua.
  \item \textbf{Cache:} embedding cache (Redis); \textbf{semantic cache} câu hỏi lặp (so cosine $>0.95$ $\rightarrow$ trả cache).
  \item \textbf{Pipeline streaming:} không chờ cả lô; file nào xong index ngay.
  \item \textbf{Serving:} vLLM batch; \emph{fast-path} SLM rẻ, escalate model lớn khi confidence thấp.
  \item \textbf{Mục tiêu đo:} throughput (file/phút), p95 latency truy vấn.
\end{itemize}

\subsection{Pha 8 --- Đánh giá}
\begin{longtable}{>{\raggedright\arraybackslash}p{4.2cm} >{\raggedright\arraybackslash}p{6.2cm} >{\raggedright\arraybackslash}p{4.2cm}}
\toprule
\textbf{Tầng} & \textbf{Metric} & \textbf{Mục tiêu} \\
\midrule
OCR & CER/WER (jiwer) & CER $<3\%$ digital, $<8\%$ scan \\
Bảng & TEDS & $>0.90$ \\
Trích xuất số & Cell accuracy & $\geq 95\%$ \\
Ánh xạ mã số & Mapping accuracy & $\geq 90\%$ \\
Đối chiếu & Catch-rate (synthetic cấy lỗi) & $\geq 90\%$ \\
Rủi ro & Backtest công ty đã vỡ/scandal & flag đúng $\geq 1$ ca \\
Agent & \% kết luận số có trích dẫn & $100\%$ (không bịa) \\
\bottomrule
\end{longtable}

\subsection{Pha 9 --- Vận hành}
Langfuse (trace agent/token-cost) + Prometheus/Grafana (hệ thống); audit log mọi kết luận; phân quyền duyệt; không dùng PII thật.

% =====================================================================
\section{Phụ lục --- Bảng tổng hợp ngưỡng \& tham số}
\begin{longtable}{>{\raggedright\arraybackslash}p{6.5cm} >{\raggedright\arraybackslash}p{8.3cm}}
\toprule
\textbf{Tham số} & \textbf{Giá trị mặc định} \\
\midrule
\endhead
Phát hiện digital (n\_chars/trang) & $\geq 50$ \& phủ $>5\%$ \\
Render scan & 300 DPI \\
Fuzzy match (RapidFuzz token\_sort) & $\geq 88$ \\
Embedding accept (cosine) & $> 0.75$ \\
Semantic cache hit (cosine) & $> 0.95$ \\
LoRA (Qwen2-VL-2B) & r=16, alpha=32, dropout=0.05, lr=1e-4 \\
LoRA (Donut decoder) & r=8..16 \\
Conformal coverage $1-\alpha$ & 0.90 \\
Calibration set & 200--500 mẫu có nhãn người \\
Agent verify khi & confidence $<0.7$ hoặc chạm quy định \\
Agent abstain khi & confidence $<0.5$ hoặc không có bằng chứng \\
Finetune trigger & coverage $<0.90$ \emph{hoặc} PSI $>0.2$ \emph{hoặc} $\geq 50$ correction \\
Oversample correction & $\times 3$ \\
Promote gate ($\varepsilon$) & accuracy không tụt quá $0.5\%$ \\
Giữ adapter & $\geq 3$ bản gần nhất \\
\bottomrule
\end{longtable}

\vspace{0.4cm}
\begin{center}\color{fsgray}\small --- Hết bản spec chi tiết FinSight ---\end{center}

\end{document}
